% common/macros.tex — dos macros de ayuda, compartidos por todos los documentos
% del curso (avance, final).

% ---------------------------------------------------------------------------
% Modo FINAL: `make es FINAL=1` crea el archivo .final-mode y las cajas de ayuda
% desaparecen. El flag lo busca ESTE archivo con \IfFileExists, así que es
% relativo al directorio desde donde compila latexmk (la carpeta del documento).
% ---------------------------------------------------------------------------
\newif\ifdocFinal \docFinalfalse
\IfFileExists{.final-mode}{\docFinaltrue}{}

% ---------------------------------------------------------------------------
% \limite{N}{qué pide la consigna aquí}
%
% Caja al inicio de cada sección con el presupuesto de palabras y el texto
% literal de la consigna (intro.md en el avance, formulario oficial en el
% final). Visible en borrador, invisible en FINAL.
% N = 0 significa "sin límite declarado".
% ---------------------------------------------------------------------------
\newcommand{\limite}[2]{%
  \ifdocFinal\else
    \begin{tcolorbox}[
      colback=blue!3, colframe=blue!35, boxrule=0.4pt,
      left=6pt, right=6pt, top=4pt, bottom=4pt,
      before skip=6pt, after skip=10pt]
      \footnotesize
      \ifnum#1>0\textbf{Máximo #1 palabras.}\quad\fi
      #2
    \end{tcolorbox}
  \fi
}

% ---------------------------------------------------------------------------
% \nota{...} — recordatorio interno de qué falta o de cómo se conecta con el
% resto del proyecto. Invisible en FINAL.
% ---------------------------------------------------------------------------
\newcommand{\nota}[1]{%
  \ifdocFinal\else
    {\color{orange!70!black}\footnotesize\textsf{[nota: #1]}\par\medskip}
  \fi
}

% ---------------------------------------------------------------------------
% \TODO{...} — dato que todavía no se tiene. Nunca se inventa un dato: se
% marca. En FINAL sigue VISIBLE y en rojo, a propósito: un documento no se
% entrega con campos sin llenar, y el modo final tiene que gritarlo.
% ---------------------------------------------------------------------------
\newcommand{\TODO}[1]{{\color{red}\textbf{[POR DEFINIR: #1]}}}
